\section{Discussion and Bibliography}
\label{sec:2.3}

In this chapter, we took as our starting point a random number generator to produce uniform
draws. The study of random number generators is a rich subject in and on itself; we refer to
\cite{83,127,136} for in-depth surveys of existing techniques and philosophical discussions on the
use of deterministic algorithms for random simulation.

Transformation methods and accept/reject methods for non-uniform sampling are reviewed
in many textbooks, see e.g.\ \cite{56,193}. The first algorithmic use of rejection sampling is often
credited to Von Neumann \cite{232} although the idea is already underpinning Buffon's needle experiment. The paper \cite{232} also considers inversion from a uniform proposal distribution. Another
early example of a transformation method is the Box-Muller algorithm \cite{33} to sample from a
normal distribution, which you will derive in Exercise~2.3. The KR rearrangement was introduced in \cite{126,204}. We refer to \cite{7} for further background on disintegration of measures. More
recently, \cite{157} investigates the use of transport maps for sampling. Triangular transport maps
also form the building blocks of many complex architectures for \textit{normalizing flows} in machine
learning \cite{181}.

We refer to \cite{227} for an accessible introduction to ABC. The main ideas behind ABC can
be traced back at least to \cite{205}, but it was not until the late 90s that ABC methods received
their name and gained popularity through a series of influential works, including \cite{224,188,18}.
The ABC rejection sampler in Algorithm~\ref{alg:2.3} was introduced in \cite{188}. Implementations of ABC
using Markov chain Monte Carlo algorithms and particle filters were introduced in \cite{155,214}.
ABC originates in the population genetics literature, where models are inherently complex and
likelihoods are usually prohibitively costly to evaluate or impossible to compute. ABC techniques are also widely used in genetics \cite{235,213,75}, epidemiology \cite{223}, human demographic
history \cite{104}, phylogeography \cite{19}, and many other applications. From a theoretical viewpoint,
an influential paper is \cite{238}, which shows that ABC gives exact inference for a wrong model.
ABC methods for Bayesian model choice have been extensively investigated, see for instance
\cite{195} and references therein. An active area of research relies on conditional density estimation
for likelihood-free inference, see \cite{180} and also \cite{182}. These techniques have the advantage of
not relying on heuristically chosen tolerance levels.
